\section[Práctica con racionales y división]{Ejercicios con productos, números racionales y división}

En los siguientes ejercicios utilizaremos conjuntamente las ideas desarrolladas en la segunda parte del capítulo. Conviene distinguir siempre entre el número representado y la forma particular utilizada para escribirlo. Una fracción, una división y una escritura decimal pueden, en ciertos casos, representar exactamente la misma cantidad.

\subsection*{Producto de números enteros}

Antes de evaluar cada producto, indique si el resultado será positivo, negativo o cero. Luego encuentre su valor.

\begin{enumerate}
\item \((-3)\cdot4\)
\item \(5\cdot(-2)\)
\item \((-6)\cdot(-3)\)
\item \((-7)\cdot1\)
\item \((-8)\cdot0\)
\item \((-1)\cdot12\)
\item \((-4)\cdot(-5)\)
\item \(9\cdot(-6)\)
\item \((-2)\cdot3\cdot4\)
\item \((-2)\cdot(-3)\cdot5\)
\item \((-1)\cdot(-1)\cdot(-1)\)
\item \((-4)\cdot2\cdot(-3)\)
\end{enumerate}

\subsection*{Completar productos con enteros}

Complete los espacios con un número entero que haga verdadera cada igualdad.

\begin{enumerate}
\item \((-3)\cdot\square=-12\)
\item \(5\cdot\square=-20\)
\item \((-6)\cdot\square=18\)
\item \(\square\cdot(-4)=20\)
\item \((-7)\cdot\square=0\)
\item \((-1)\cdot\square=9\)
\item \((-3)\cdot(-5)=3\cdot\square\)
\item \(4\cdot(-6)=(-4)\cdot\square\)
\end{enumerate}

\subsection*{Representaciones fraccionarias equivalentes}

Complete los espacios de manera que las fracciones representen el mismo número racional.

\begin{enumerate}
\item \(\dfrac12=\dfrac{\square}{4}\)
\item \(\dfrac23=\dfrac{\square}{12}\)
\item \(\dfrac35=\dfrac{12}{\square}\)
\item \(\dfrac47=\dfrac{20}{\square}\)
\item \(\dfrac{-2}{3}=\dfrac{-8}{\square}\)
\item \(\dfrac56=\dfrac{\square}{30}\)
\item \(\dfrac{7}{10}=\dfrac{21}{\square}\)
\item \(\dfrac{-3}{4}=\dfrac{\square}{20}\)
\item \(3=\dfrac{\square}{5}\)
\item \(-2=\dfrac{\square}{7}\)
\end{enumerate}

Ahora escriba tres representaciones fraccionarias diferentes para cada uno de los siguientes números:
\[
\frac12,\qquad \frac34,\qquad -\frac25,\qquad 2.
\]

\subsection*{Reconocer números racionales}

Responda en cada caso y justifique brevemente utilizando la definición de número racional.

\begin{enumerate}
\item ¿Por qué \(5\) es un número racional aunque esté escrito como un entero?
\item ¿Por qué \(-3\) es un número racional?
\item ¿La expresión \(\frac{4}{0}\) representa un número racional? Explique.
\item ¿La expresión \(\frac{0}{7}\) representa un número racional? ¿Qué número representa?
\item Escriba \(8\) como una fracción cuyo denominador sea \(3\).
\item Escriba \(-5\) como una fracción cuyo denominador sea \(4\).
\end{enumerate}

\subsection*{Producto de números racionales}

Encuentre una fracción que represente el valor de cada producto. Puede reemplazar la fracción obtenida por una representación equivalente más sencilla cuando sea posible.

\begin{enumerate}
\item \(\dfrac12\cdot\dfrac34\)
\item \(\dfrac23\cdot\dfrac57\)
\item \(\dfrac35\cdot\dfrac{10}{9}\)
\item \(\dfrac47\cdot\dfrac{14}{3}\)
\item \(\dfrac{-2}{3}\cdot\dfrac54\)
\item \(\dfrac34\cdot\dfrac{-8}{5}\)
\item \(\dfrac{-5}{6}\cdot\dfrac{-3}{10}\)
\item \(2\cdot\dfrac37\)
\item \((-4)\cdot\dfrac58\)
\item \(\dfrac79\cdot0\)
\item \(\dfrac{11}{5}\cdot1\)
\item \(\dfrac23\cdot\dfrac34\cdot\dfrac56\)
\end{enumerate}

\subsection*{Inversos multiplicativos}

Escriba el inverso multiplicativo de cada número, cuando exista.

\begin{enumerate}
\item \(\dfrac23\)
\item \(\dfrac57\)
\item \(-\dfrac34\)
\item \(4\)
\item \(-6\)
\item \(\dfrac{1}{9}\)
\item \(-1\)
\item \(1\)
\item \(\dfrac{-7}{2}\)
\item \(0\)
\end{enumerate}

En cada caso en que exista un inverso, compruebe su respuesta multiplicando ambos números y verificando que el resultado sea \(1\). En el caso de \(0\), explique por qué esa comprobación no puede realizarse.

\subsection*{Reescribir divisiones como productos}

Reescriba cada división como un producto por el inverso multiplicativo del divisor. No es necesario evaluarla todavía.

\begin{enumerate}
\item \(8\div2\)
\item \(5\div3\)
\item \((-6)\div2\)
\item \(7\div(-4)\)
\item \(\dfrac34\div2\)
\item \(\dfrac25\div\dfrac37\)
\item \(\dfrac{-4}{9}\div\dfrac23\)
\item \(3\div\dfrac58\)
\item \((-2)\div\dfrac{-7}{3}\)
\item \(\dfrac56\div(-5)\)
\end{enumerate}

Por ejemplo,
\[
\frac34\div\frac25
=
\frac34\cdot\frac52.
\]
La segunda expresión no representa una operación diferente: representa la misma división escrita mediante la definición \(x\div y=x\cdot y^{-1}\).

\subsection*{Evaluar divisiones}

Encuentre una fracción que represente el valor de cada expresión.

\begin{enumerate}
\item \(8\div2\)
\item \(15\div3\)
\item \(7\div4\)
\item \(1\div2\)
\item \((-12)\div3\)
\item \(18\div(-6)\)
\item \((-20)\div(-5)\)
\item \(\dfrac34\div\dfrac25\)
\item \(\dfrac56\div\dfrac23\)
\item \(\dfrac{-2}{7}\div\dfrac35\)
\item \(\dfrac49\div\dfrac{-8}{3}\)
\item \(2\div\dfrac37\)
\item \((-5)\div\dfrac{10}{3}\)
\item \(\dfrac78\div2\)
\item \(\dfrac32\div\dfrac94\)
\end{enumerate}

\subsection*{Completar igualdades con productos y divisiones}

Complete cada espacio con un número racional que haga verdadera la igualdad.

\begin{enumerate}
\item \(2\cdot\square=1\)
\item \(5\cdot\square=1\)
\item \((-4)\cdot\square=1\)
\item \(\dfrac37\cdot\square=1\)
\item \(\dfrac{-5}{2}\cdot\square=1\)
\item \(\square\div2=3\)
\item \(8\div\square=4\)
\item \(\dfrac34\div\square=\dfrac32\)
\item \(\square\div\dfrac25=\dfrac32\)
\item \(\dfrac56\cdot\square=\dfrac{5}{12}\)
\end{enumerate}

Observe que algunos de estos ejercicios pueden interpretarse como ecuaciones multiplicativas. Utilice la idea de inverso multiplicativo para justificar sus respuestas cuando resulte conveniente.

\subsection*{Restas sucesivas}

Utilice restas sucesivas para evaluar las siguientes divisiones. En cada caso, indique qué representa la cantidad de restas realizadas.

\begin{enumerate}
\item \(8\div2\)
\item \(12\div3\)
\item \(20\div4\)
\item \(18\div6\)
\item \(35\div5\)
\item \(42\div7\)
\end{enumerate}

Luego responda: ¿por qué este procedimiento no resulta suficiente para explicar directamente una división como \(1\div2\)?

\subsection*{Fracciones y escritura decimal}

Escriba cada fracción mediante una escritura decimal equivalente.

\begin{enumerate}
\item \(\dfrac{1}{10}\)
\item \(\dfrac{3}{10}\)
\item \(\dfrac{7}{10}\)
\item \(\dfrac{25}{100}\)
\item \(\dfrac{3}{4}\)
\item \(\dfrac12\)
\item \(\dfrac{5}{2}\)
\item \(\dfrac{7}{4}\)
\item \(\dfrac{15}{20}\)
\item \(\dfrac{125}{100}\)
\end{enumerate}

Ahora realice el proceso contrario: escriba cada número decimal como una fracción de enteros con denominador distinto de cero.

\begin{enumerate}
\item \(0{,}1\)
\item \(0{,}5\)
\item \(0{,}25\)
\item \(0{,}7\)
\item \(0{,}75\)
\item \(1{,}5\)
\item \(1{,}25\)
\item \(2{,}5\)
\item \(0{,}04\)
\item \(3{,}125\)
\end{enumerate}

\subsection*{Distintas representaciones del mismo número}

Complete los espacios de manera que todas las expresiones de cada fila representen el mismo número racional.

\begin{enumerate}
\item \(1\div2=\dfrac{\square}{2}=0{,}5\)
\item \(1\div4=\dfrac14=\square\)
\item \(3\div4=\dfrac34=\square\)
\item \(5\div2=\dfrac52=\square\)
\item \(7\div4=\dfrac74=\square\)
\item \(\dfrac{25}{100}=\dfrac{\square}{4}=0{,}25\)
\item \(0{,}5=\dfrac{5}{\square}=\dfrac12\)
\item \(1{,}25=\dfrac{125}{\square}=\dfrac54\)
\end{enumerate}

En estos ejercicios no aparecen números diferentes en cada miembro. Se muestran escrituras distintas de una misma cantidad.

\subsection*{Algoritmo de la división}

Utilice el algoritmo de la división o una descomposición equivalente basada en múltiplos del divisor. Muestre los pasos principales.

\begin{enumerate}
\item \(84\div7\)
\item \(96\div8\)
\item \(144\div12\)
\item \(156\div12\)
\item \(225\div15\)
\item \(324\div18\)
\item \(1\div2\)
\item \(3\div4\)
\item \(7\div4\)
\item \(9\div8\)
\end{enumerate}

En los casos en que el cociente no sea entero, continúe trabajando con décimos, centésimos o milésimos según resulte necesario.

\subsection*{Pertenencia a conjuntos numéricos}

Indique con una cruz (\(\times\)) a cuáles de los conjuntos \(\mathbb N\), \(\mathbb Z\) y \(\mathbb Q\) pertenece cada número. Un mismo número puede pertenecer a más de un conjunto.

\begin{center}
\begin{tabular}{c|c|c|c}
  Número & Natural (\(\mathbb{N}\)) & Entero (\(\mathbb{Z}\)) & Racional (\(\mathbb{Q}\)) \\ \hline
\(10\) & & & \\ \hline
\(-7\) & & & \\ \hline
\(0\) & & & \\ \hline
\(3/4\) & & & \\ \hline
\(-5/2\) & & & \\ \hline
\(5{,}0\) & & & \\ \hline
\(0{,}5\) & & & \\ \hline
\(12/3\) & & & \\ \hline
\(-8/4\) & & & \\ \hline
\(0/9\) & & & \\ \hline
\(1{,}25\) & & & \\ \hline
\(-3{,}0\) & & &  
\end{tabular}
\end{center}

En los casos en que la escritura pueda ocultar el conjunto más pequeño al que pertenece el número, evalúe o reescriba primero la expresión. Por ejemplo, una fracción puede representar un entero.

\subsection*{Preguntas conceptuales}

Responda con una explicación breve. Cuando sea posible, acompañe su respuesta con una igualdad o con un ejemplo.

\begin{enumerate}
\item ¿Qué diferencia hay entre los factores de un producto y los sumandos de una suma?
\item ¿Por qué \(a\cdot0=0\) para cualquier número \(a\)?
\item ¿Por qué el producto por \(1\) no modifica el número representado?
\item ¿Por qué \((-a)\cdot(-b)=a\cdot b\)? Relacione su respuesta con la idea de opuesto.
\item ¿Por qué una misma cantidad racional puede escribirse mediante fracciones diferentes?
\item ¿Por qué el denominador de una fracción racional no puede ser cero?
\item ¿Por qué \(0\) no posee inverso multiplicativo?
\item ¿Qué significa afirmar que dividir por un número no nulo equivale a multiplicar por su inverso?
\item Explique por qué \(1\div2\), \(\frac12\) y \(0{,}5\) representan el mismo número.
\item ¿Por qué las restas sucesivas son una estrategia útil para algunas divisiones, pero no una definición general de la división?
\item ¿La escritura decimal crea un conjunto numérico nuevo? Justifique con un ejemplo del capítulo.
\item Explique con sus palabras la relación \(\mathbb N\subseteq\mathbb Z\subseteq\mathbb Q\).
\end{enumerate}

\subsection*{Construir expresiones}

Construya una expresión que cumpla cada condición. En muchos casos existen varias respuestas correctas.

\begin{enumerate}
\item Un producto de dos naturales que represente \(24\).
\item Dos productos diferentes que representen el mismo número natural.
\item Un producto de dos enteros negativos que represente \(30\).
\item Un producto de dos enteros cuyo resultado sea \(-24\).
\item Una fracción diferente de \(\frac12\) que represente el mismo número.
\item Dos fracciones diferentes que representen \(-\frac34\).
\item Un producto de dos racionales que represente \(1\), sin utilizar \(1\) como factor.
\item Una división cuyo resultado sea \(\frac12\).
\item Una división entre dos enteros cuyo resultado no sea entero.
\item Tres escrituras diferentes que representen \(0{,}5\).
\item Un número que pertenezca a \(\mathbb Q\) pero no a \(\mathbb Z\).
\item Un número que pertenezca a \(\mathbb Z\) pero no a \(\mathbb N\).
\end{enumerate}
